\chapter{Conclusion}
\label{sec:Conclusion}

The velocity field reconstructed in this thesis reproduces the bulk motion of the Taurus complex
and, once that motion is subtracted, separates its subgroups by their internal velocities. The
age--velocity gradient already visible in the stellar data persists in the field and across most
runs, and the divergence and vorticity maps add a qualitative picture that does not depend on the
chosen reference frame. What the field supports are statements about which parts of the complex move differently from each
other, rather than measurements of individual velocities: if the bulk motion is off, every
vector of the field is off by the same amount, so the error cancels when two regions are compared
but remains in any velocity value quoted on its own.

The caveats are as important as the result. The mean motion, and with it the internal reference
frame, varies by several $\mathrm{km\,s^{-1}}$ between runs, which enters every absolute magnitude and
direction in that frame, while a comparison between two regions of the same run is unaffected by
that shift. Small-scale peaks appear
throughout the region but at different positions in different runs, which marks them as
consequences of individual stars rather than properties of the complex. The reconstruction also
inherits the geometry of the region: $U$ points almost along the line of sight and is therefore the
component that the scarce radial velocities have to constrain, which makes it the least certain
part of the field and the source of its most conspicuous artifacts. Underlying all of this is the
central assumption that the stars follow a smooth field. They do so only approximately, and the model
offers no place for the difference: a peculiar motion ends up either in the field or in the
residuals. Granting the stars a common intrinsic dispersion repairs the
residual statistics, but the value obtained this way is a fit parameter rather than a measurement,
since the same data constrain it and the amplitude of the field.

Most of these limitations trace back to one source: too few stars, and radial velocities for only a
fraction of them. With 557 members on a grid of about $26\,000$ cells and 123 usable radial velocities,
the data cannot constrain the field, smooth the noise and absorb outliers at the same time. This is
not a shortcoming of the method, which has been applied successfully to the far richer sample of
the Scorpius-Centaurus association \citep{Hutschenreuter2026}, but of what the region currently
offers. Gaia~DR4 will improve
exactly this: radial velocities derived from the Radial Velocity Spectrometer are announced for
about $312$ million sources, and for the first time also at the level of the individual transits
\citep{GaiaDR4}. The latter would turn the stellar variability that this work had to absorb into an
error term into a measurable quantity. Repeating the reconstruction on that basis is the most direct
way to test which of the features described here survive. The model itself also
leaves room to grow: a separate scatter term for the radial velocities, which the quality statistics of this
work already call for.
